\chapter{Related Work}
\label{chap:related-work}

\section{Federated Learning}

FedAvg~\cite{mcmahan2017fedavg} is the foundational FL algorithm and the
training protocol assumed throughout this project: clients run several
steps of local SGD on their own data and the server aggregates the
resulting updates by a weighted average (Equation~\ref{eq:fedavg} in
Chapter~\ref{chap:problem-statement}). Its privacy argument is informal -
clients never transmit raw data, only model updates - and is exactly the
assumption the gradient inversion literature below sets out to test.

\section{Optimization-Based Gradient Inversion}

\textbf{Deep Leakage from Gradients (DLG)}~\cite{zhu2019dlg} was the first
to show that a single gradient can leak the private image and label that
produced it. DLG treats reconstruction as an optimization problem: a dummy
input and label are updated by gradient descent so that their gradient
matches the leaked one under an $L_2$ distance, using second-order
differentiation (``double backpropagation'') to backpropagate through the
gradient-matching loss itself. Joint optimization of the dummy label
alongside the image makes convergence slow and occasionally unstable.

\textbf{iDLG}~\cite{zhao2020idlg} removes that instability by observing
that, for cross-entropy loss on a one-hot label, the sign pattern of the
last layer's gradient identifies the ground-truth class in closed form -
only the dummy image still needs to be optimized. This simplification
makes convergence markedly more reliable.

\textbf{Inverting Gradients}~\cite{geiping2020inverting} argues that most
of the information in a gradient is carried by its \emph{direction} rather
than its magnitude, and replaces the $L_2$ distance with a
cosine-similarity cost. It further adds a total-variation image prior and
optimizes with a sign-based variant of Adam. These changes make the attack
scale from small networks up to ImageNet-scale ResNet architectures,
showing that gradient leakage is a practical threat even for deep
convolutional networks trained with non-trivial batch sizes.

\section{Closed-Form Gradient Inversion}
\label{sec:related-closed-form}

\textbf{R-GAP (Recursive Gradient Attack on Privacy)}~\cite{zhu2021rgap}
reformulates gradient leakage as a linear algebra problem rather than a
search problem. For a feed-forward network with bias-free layers and
invertible activations, the weight-gradient equation at each layer relates
a known quantity (the leaked per-layer weight gradient) to the layer's
input, which can be solved for directly; recursing this from the output
layer back to the input recovers the private sample with no optimization
loop. R-GAP additionally proposes a \emph{rank analysis} that predicts,
from the architecture alone, whether a given layer's input is uniquely
recoverable, without running the attack at all.

\section{Defenses Against Gradient Inversion}

\textbf{DP-SGD}~\cite{abadi2016dpsgd} was originally proposed to give
differential-privacy guarantees during training, by clipping each
per-example gradient to a norm $C$ and adding Gaussian noise with
multiplier $\sigma$. The same clip-and-noise procedure is widely reused, as
in this project, as an ad hoc empirical defense against gradient inversion:
since it is applied directly to the gradient the attacker observes, it
perturbs the attacker's input rather than any property of the model.

\section{Positioning of This Project}

This project contributes an empirical comparison, under matched conditions,
of optimization-based attacks (DLG/iDLG, and separately Inverting
Gradients) and a closed-form attack (R-GAP) across comparable victim
architectures, together with an evaluation of how all three degrade under
the same DP-SGD defense budget.
